\section*{Chapitre \ref{ch:b3:planetary-interiors-surfaces} --- Intérieurs et surfaces planétaires}

\begin{solution}{exo:b3:planetary-interiors-surfaces:1}
$\tfrac52C\bar\rho = 2.5\times0.35\times5429 = 4750$; $x^2 = (4750 -
3300)/(5429 - 3300) = 0.681$, $x = 0.825$, soit un noyau de $0.825\times2440 =
\qty{2010}{km}$; $\rho_c = 3300 + 2129/0.825^3 = \qty{7090}{kg.m^{-3}}$.
\end{solution}

\begin{solution}{exo:b3:planetary-interiors-surfaces:2}
$m = 3500\times\tfrac43\pi(25)^3 = \qty{2.29e8}{kg}$; $E = \tfrac12 m v^2 =
\tfrac12\times\num{2.29e8}\times(\num{2.0e4})^2 = \qty{4.6e16}{J}$, ou
$\num{4.6e16}/\num{4.2e15} = \qty{11}{Mt}$. La règle donne un cratère de
$20\times\qty{50}{m} = \qty{1}{km}$.
\end{solution}

\begin{solution}{exo:b3:planetary-interiors-surfaces:3}
$N_1 = 120/\num{100000} = \qty{1.2e-3}{km^{-2}}$, erreur relative $1/\sqrt{120}
= \qty{9}{\percent}$. Sur la partie linéaire, $T = N_1/c = \num{1.2e-3}/\num{4.875e-4}
= \qty{2.46}{Ga}$ (le terme exponentiel, $\num{8e-7}$, est négligeable), et
$\sigma_T = T/\sqrt{120} = \qty{0.22}{Ga}$.
\end{solution}

\begin{solution}{exo:b3:planetary-interiors-surfaces:4}
Aire $4\pi(\num{1.8215e6})^2 = \qty{4.17e13}{m^2}$~; puissance $2.5\times
\num{4.17e13} = \qty{1.04e14}{W} = \qty{104}{TW}$, soit 2.2 fois le total terrestre
pour une surface douze fois plus petite, et $2.5/0.092 = 27$ fois son
flux.
\end{solution}

\begin{solution}{exo:b3:planetary-interiors-surfaces:5}
(a)~$\bar\rho = \num{6.417e23}/(\tfrac43\pi(\num{3.3895e6})^3) =
\qty{3934}{kg.m^{-3}}$; $g = GM/R^2 = \qty{3.73}{m.s^{-2}}$; $\rho gR =
\qty{4.97e10}{Pa}$. (b)~$19\mu/(2\rho gR) = 19\times\num{6.7e10}/\num{9.94e10}
= 12.8$, $k_2 = 1.5/13.8 = 0.109$. (c)~La valeur mesurée est supérieure de
2.6 écarts-types~: Mars donne une valeur plus élevée que celle d'une planète
solide dotée d'une rigidité mantellique, comme il le ferait avec un noyau liquide interne.
\end{solution}

\begin{solution}{exo:b3:planetary-interiors-surfaces:6}
(a)~Lune~: $4\times9.8/1.62 = \qty{24}{km}$ dans la roche cristalline,
\qty{12}{km} dans les sédiments~; Mars et Mercure~: $4\times9.8/3.7 =
\qty{10.6}{km}$ et \qty{5.3}{km}. Les \qty{20}{km} de la Lune se situent entre
ses deux prévisions~: la règle tient à quelques dizaines de pour cent. (b)~$Y =
\rho gD/3 = 2600\times9.8\times4000/3 = \qty{34}{MPa}$, soit la résistance
d'une roche fracturée plutôt que d'un cristal intact.
\end{solution}

\begin{solution}{exo:b3:planetary-interiors-surfaces:7}
Mars~: $\num{3.22e-14}(e^{6.93T} - 1) + \num{4.875e-4}\,T = \num{2.0e-3}$~;
l'estimation linéaire $4.1$~Ga est beaucoup trop ancienne, le terme exponentiel étant grand
là-bas~; la dichotomie donne $T = \qty{3.34}{Ga}$ (terme exponentiel
$\num{3.6e-4}$, terme linéaire $\num{1.63e-3}$). La fonction lunaire donne
\qty{2.39}{Ga}. Chaque fonction est calibrée pour un seul corps~: le taux
d'impacts et la taille des cratères qu'un projectile donné creuse diffèrent
entre la Lune et Mars.
\end{solution}

\begin{solution}{exo:b3:planetary-interiors-surfaces:8}
(a)~$3.551181/1.769138 = 2.0073$ et $7.154553/3.551181 = 2.0147$, proche de
$2$. (b)~$0.565247 - 0.844789 + 0.279542 = 0.000000$ par jour, à
$\num{1e-7}$ près~: les trois moyens mouvements sont verrouillés, une résonance qui maintient
les positions relatives des lunes à chaque conjonction. (c)~Les attractions se répètent
aux mêmes points de l'orbite, de sorte qu'elles s'additionnent au lieu de se compenser~;
elles alimentent l'excentricité aussi vite que les marées l'amortissent.
\end{solution}

\begin{solution}{exo:b3:planetary-interiors-surfaces:9}
$n = 2\pi/(1.769138\times86400) = \qty{4.111e-5}{s^{-1}}$, $nR =
\qty{74.9}{m.s^{-1}}$, $(nR)^5/G = \qty{3.53e19}{W}$;
$\tfrac{21}{2}e^2(nR)^5/G = 10.5\times\num{1.6e-5}\times\num{3.53e19} =
\qty{5.9e15}{W}$ par unité de $k_2/Q$. Puis $k_2/Q =
\num{1.04e14}/\num{5.9e15} = 0.018$.
\end{solution}

\begin{solution}{exo:b3:planetary-interiors-surfaces:10}
(a)~$x^2 = (\tfrac52C\bar\rho - \rho_m)/(\bar\rho - \rho_m)$ est linéaire en
$C$, de pente $\tfrac52\bar\rho/(\bar\rho - \rho_m)$. (b)~La pente vaut
$2.5\times3340/52 = 161$~: $x^2$ varie de $0.161$ pour $0.001$ de $C$.
Pour $C = 0.395$, $x^2 = 0.036 + 0.161 = 0.197$, $x = 0.44$, $\rho_c =
\qty{3880}{kg.m^{-3}}$~; pour $C = 0.393$, $x^2 < 0$~: pas de noyau du tout.
(c)~Au voisinage de $C = \tfrac25$, le moment d'inertie ressent à peine un petit noyau~:
son troisième chiffre significatif tranche entre un noyau grand et léger et l'absence de noyau.
Les ondes sismiques, les marées et le champ magnétique sont nécessaires.
\end{solution}

\begin{solution}{exo:b3:planetary-interiors-surfaces:11}
(a)~$\sigma_T = N_1/(10\,N_1')$~: à \qty{1.0}{Ga}, $N_1 \approx N_1' \approx
\num{4.875e-4}$, $\sigma_T = \qty{0.10}{Ga}$~; à \qty{3.8}{Ga}, $N_1 =
\num{1.07e-2}$ et $N_1' = \num{6.15e-2}$, $\sigma_T = \qty{0.017}{Ga}$,
six fois plus petit. (b)~Avec $N_1 = cT$, $N_1' = c$ et $\sigma_T =
cT/(\sqrt n\,c) = T/\sqrt n$. (c)~La fonction de chronologie elle-même est
incertaine~: elle repose sur quelques sites datés de la Lune et sur la mise à l'échelle
de la Lune à Mars, des erreurs qu'aucun comptage ne réduit.
\end{solution}

\begin{solution}{exo:b3:planetary-interiors-surfaces:12}
(a)~Le rapport vaut $(n_ER_E/n_IR_I)^5(e_E/e_I)^2 = (0.4982\times0.8569)^5
\times(0.009/0.004)^2 = 0.0142\times5.06 = 0.072$~: soit environ
\qty{7.5e12}{W}, ou $\num{7.5e12}/(4\pi(\num{1.5608e6})^2) =
\qty{0.24}{W.m^{-2}}$. (b)~Un océan découple la coquille de glace de
l'intérieur~: la coquille se déforme bien davantage qu'une lune solide, et sa glace
se dissipe autrement que la roche chaude d'Io, de sorte que $k_2$ et $Q$ changent tous deux.
\end{solution}

\begin{solution}{pb:b3:planetary-interiors-surfaces:1}
\textbf{1.} $\bar\rho = \num{6.4169e23}/(\tfrac43\pi(\num{3.3895e6})^3) =
\qty{3934}{kg.m^{-3}}$.

\textbf{2.} $\tfrac52C\bar\rho = 3600$; $x^2 = (3600 - 3500)/(3934 - 3500) =
0.230$, $x = 0.479$, un noyau de \qty{1620}{km}~; $\rho_c = 3500 +
434/0.479^3 = \qty{7450}{kg.m^{-3}}$.

\textbf{3.} $x = 1830/3389.5 = 0.540$; $\rho_m = (3600 - 0.2915\times
3934)/(1 - 0.2915) = \qty{3460}{kg.m^{-3}}$, $\rho_c = 3462 +
472/0.1574 = \qty{6460}{kg.m^{-3}}$.

\textbf{4.} C'est un tiers de moins. La pression au cœur de Mars est bien
plus faible que sur la Terre, si bien que son métal est moins comprimé~; et
un noyau aussi léger suggère des éléments bien plus légers mêlés au fer.

\textbf{5.} Le \omterm{def:b3:planetary-interiors-surfaces:love-number}{nombre de Love} (2003) est bien antérieur à InSight (2021)~: il
montrait déjà que le noyau externe au moins est liquide, et plaçait son rayon
entre 1520 et \qty{1840}{km}, un intervalle qui contient la valeur
sismologique.

\textbf{6.} $1 - (1 - 20/3389.5)^3 = \qty{1.8}{\percent}$ à $1 - (1 -
37/3389.5)^3 = \qty{3.2}{\percent}$ du volume, et moins de la masse~:
plus léger que le manteau, il abaisse un peu la densité effective du manteau,
un effet plus petit que l'incertitude du modèle.

\textbf{7.} À partir de \qty{1}{km}~: 54~; à partir de 1.41~: 16~; à partir de 2~: 10~; à partir de 2.83~: 5~;
à partir de 4, 5.66, 8 et 11.3~: 2~; à partir de 16~: 1.

\textbf{8.} $N = 54/80000 = \qty{6.75e-4}{km^{-2}}$, $\sigma = \sqrt{54}/80000
= \qty{0.92e-4}{km^{-2}}$ (\qty{14}{\percent}).

\textbf{9.} Attendu~: $54\times0.355 = 19$ cratères de \qty{1.41}{km} et
$54\times0.147 = 7.9$ de \qty{2}{km}~; comptés 16 et 10, à moins d'un écart
type ($\sqrt{19} = 4.4$, $\sqrt{7.9} = 2.8$)~: cohérent.

\textbf{10.} Non~: $0.072/\num{6.75e-4} = 107$ fois en dessous. Le comptage
enregistre le temps.

\textbf{11.} Des intervalles de même largeur en $\log D$ couvrent uniformément la plage
et conservent des effectifs comparables dans les classes d'une distribution raide.

\textbf{12.} Ce sont des cratères secondaires, produits par des blocs éjectés
d'un grand cratère primaire~: ce ne sont pas des impacts indépendants venus de
l'espace, ils ajouteraient des cratères que la fonction de chronologie ne compte pas.

\textbf{13.} $T = \num{6.75e-4}/\num{4.875e-4} = \qty{1.385}{Ga}$~; le
terme exponentiel, $\num{3.22e-14}(e^{9.6} - 1) = \num{4.7e-10}$, vaut un
millionième de $N_1$.

\textbf{14.} $\sigma_T = T/\sqrt{54} = \qty{0.19}{Ga}$~: $1.38 \pm 0.19$~Ga.

\textbf{15.} $\num{6.75e-4}/\num{8.38e-4} = \qty{0.81}{Ga}$, plus jeune de
près de moitié.

\textbf{16.} $31 + 12 + 7 + 2 + 0 + 2 + 2 = 56$ cratères~; $N(\qty{2}{km}) =
56/120000 = \qty{4.67e-4}{km^{-2}}$~; $N_1 = \num{4.67e-4}/0.1473 =
\qty{3.17e-3}{km^{-2}}$. Dichotomie sur $\num{3.22e-14}(e^{6.93T} - 1) +
\num{4.875e-4}\,T = \num{3.17e-3}$ donne $T = \qty{3.54}{Ga}$ (terme
exponentiel \num{1.45e-3}, terme linéaire \num{1.73e-3}).

\textbf{17.} $N_1' = \num{3.22e-14}\times6.93\,e^{6.93\times3.54} +
\num{4.875e-4} = \num{1.05e-2}$~; $\sigma_T = \num{3.17e-3}/(\sqrt{56}\times
\num{1.05e-2}) = \qty{0.040}{Ga}$.

\textbf{18.} Plaine A~: \qty{14}{\percent}~; plaine B~: \qty{1.1}{\percent}. Dans
la partie initiale et raide de la chronologie, une forte variation de la densité
de cratères correspond à une faible variation d'âge~; sur la partie linéaire,
l'âge est aussi incertain que le comptage.

\textbf{19.} $120000\times\num{3.17e-3}\times(1 - 0.1473) = 324$~; elle
porte $48 + 8 = 56$.

\textbf{20.} $N_1 = 56/(120000\times0.8527) = \qty{5.47e-4}{km^{-2}}$,
$T_r = \num{5.47e-4}/\num{4.875e-4} = \qty{1.12}{Ga}$.

\textbf{21.} $\sigma = 1.12/\sqrt{56} = \qty{0.15}{Ga}$.

\textbf{22.} Un cratère de \qty{2}{km} est profond de \qty{400}{m}~: la lave
avait au moins cette épaisseur. Les cratères plus grands sont plus profonds
et leurs bourrelets se dressent plus haut~; ils ont été en partie remplis et restent visibles.

\textbf{23.} $1.38 \pm 0.19$ et $1.12 \pm 0.15$~Ga diffèrent de \qty{0.26}{Ga},
soit environ un écart type combiné ($\sqrt{0.19^2 + 0.15^2} = 0.24$)~: les
comptages permettent un épisode unique, sans le prouver.

\textbf{24.} La jeune plaine de lave A a un âge modèle de
\textbf{1.4~Ga} ($1.38 \pm 0.19$~Ga).
\end{solution}
